% ============================================================================
% Fichier  : partie4/P04_C21_evaluer.tex
% Titre    : Évaluer la Sécurité
% Auteur   : MINKA MI NGUIDJOÏ Thierry Emmanuel
% Version  : 1.0 -- novembre 2026
% ============================================================================
% Base : brouillon « Livre VII » (types de tests, étapes d'un test
% d'intrusion). Sources : NIST SP 800-115 (2008) ; ISO 19011:2018 ;
% ISO/IEC 27007:2020 ; CVSS v3.1 ; loi n° 2010/012, art. 65 (2) et 69 (texte
% officiel, cités au Chapitre 11).
% ============================================================================

\chapter{Évaluer la Sécurité}
\label{chap:evaluer}
\chaptermark{Évaluer la sécurité}

\epigraph{%
    « Un test ne trouve que ce qu'il a cherché,\\
    dans le temps qu'il avait. »%
}{Principe du cours}

\niveauChapitre{Intermédiaire}{Chapitres~\ref{chap:menaces},
\ref{chap:droit}, \ref{chap:ethique} et~\ref{chap:applications}.}

\begin{boxobjectifs}
\begin{itemize}
    \item Distinguer analyse de vulnérabilités, test d'intrusion et audit.
    \item Rédiger l'autorisation et le cadre d'un test.
    \item Décrire les phases d'un test d'intrusion.
    \item Prioriser les vulnérabilités trouvées.
    \item Lire et exploiter un rapport de test.
\end{itemize}
\end{boxobjectifs}

\bigskip

% ============================================================================
\section*{Question directrice : comment savoir que les mesures fonctionnent~?}
% ============================================================================

Six mois après l'attaque, LiZbiZ a mis en place la double
authentification, segmenté son réseau et renégocié le contrat de son
prestataire. Le conseil demande~: « sommes-nous protégés~? » Seule une
évaluation peut répondre, à condition d'en choisir le type, d'en fixer le
cadre et de lire ses résultats avec prudence.

\medskip
\textbf{Plan du chapitre.}
\begin{enumerate}
    \item Trois façons d'évaluer (section~\ref{sec:c21-types}).
    \item Le cadre et l'autorisation (section~\ref{sec:c21-cadre}).
    \item Les phases d'un test d'intrusion (section~\ref{sec:c21-phases}).
    \item Prioriser les vulnérabilités (section~\ref{sec:c21-prioriser}).
    \item Le rapport (section~\ref{sec:c21-rapport}).
\end{enumerate}

\separateur

% ============================================================================
\section{Trois Façons d'Évaluer}
\label{sec:c21-types}
% ============================================================================

\begin{tableaucours}{L{3.2cm}L{4.6cm}L{4.4cm}}{%
    \tableauHeader{Évaluation} &
    \tableauHeader{Question} &
    \tableauHeader{Résultat}
}
Analyse de vulnérabilités & Quelles faiblesses connues sont présentes~? &
Une liste, à trier \\
Test d'intrusion & Un attaquant peut-il atteindre tel objectif~? & Des
chemins d'attaque démontrés \\
Audit de sécurité & Les mesures exigées existent-elles et
fonctionnent-elles~? & Des écarts à un référentiel, avec leurs preuves \\
Exercice d'équipe rouge & L'organisation détecte-t-elle et répond-elle~?
& Une mesure de la détection et de la réponse \\
\end{tableaucours}
\vspace{-0.8em}
\begin{center}{\small\itshape Tableau~21.1. Les types d'évaluation}\end{center}

Le NIST SP~800-115 décrit ces techniques et la manière de les
organiser~\cite{nist2008sp800115}~; l'ISO~19011 et l'ISO/IEC~27007
encadrent l'audit~\cite{iso19011_2018,iso27007_2020}.

% ============================================================================
\section{Le Cadre et l'Autorisation}
\label{sec:c21-cadre}
% ============================================================================

Un test non autorisé est un accès non autorisé, puni par la loi
(art.~65, al.~2, Chapitre~\ref{chap:droit}). Tout test commence donc par
une autorisation écrite.

\begin{boxprocedure}
\textbf{La note d'autorisation d'un test}
\begin{enumerate}
    \item \textbf{Signataire} habilité à autoriser pour chaque système
          (et l'hébergeur ou le fournisseur s'il y en a un).
    \item \textbf{Périmètre}~: adresses, applications, comptes de test~;
          ce qui est exclu.
    \item \textbf{Méthodes} autorisées et interdites (déni de service,
          ingénierie sociale).
    \item \textbf{Fenêtre}~: dates et heures.
    \item \textbf{Contacts} d'urgence des deux côtés, et procédure d'arrêt.
    \item \textbf{Traitement des données} rencontrées et des résultats.
\end{enumerate}
\end{boxprocedure}

% ============================================================================
\section{Les Phases d'un Test d'Intrusion}
\label{sec:c21-phases}
% ============================================================================

\begin{tableaucours}{L{3.0cm}L{9.2cm}}{%
    \tableauHeader{Phase} &
    \tableauHeader{Contenu}
}
Planification & Objectifs, règles d'engagement, autorisation \\
Découverte & Reconnaissance passive (sources ouvertes), puis active
(balayage, énumération des services) \\
Attaque & Exploitation des vulnérabilités, élévation de privilèges,
déplacement, jusqu'à l'objectif ou la fin de la fenêtre \\
Rapport & Constatations, preuves, gravité, recommandations \\
\end{tableaucours}
\vspace{-0.8em}
\begin{center}{\small\itshape Tableau~21.2. Les phases d'un test d'intrusion
(d'après le NIST SP~800-115)}\end{center}

Les tests se conduisent en \emph{boîte noire} (sans information), en
\emph{boîte grise} (avec un compte utilisateur) ou en \emph{boîte
blanche} (avec le code et la documentation). La boîte grise est souvent
le meilleur compromis pour une application.

\begin{boxscenario}{Opération MEFOU. Le test du nouveau portail de l'ANPP}
\textbf{Cadre.} Boîte grise, deux comptes de bénéficiaires de test, du 3
au 7~mai, de 18~h à 6~h, sur une copie isolée~; déni de service exclu~;
note signée par le directeur général de l'ANPP et par l'hébergeur.

\medskip
\textbf{Constatations} (rapport du testeur, preuves en annexe)~: le
contrôle d'accès par objet est désormais vérifié (le défaut décrit au
Chapitre~\ref{chap:applications} est corrigé)~; une bibliothèque de
génération de documents présente une CVE critique exploitable à
distance~; les messages d'erreur révèlent la version du serveur.
\end{boxscenario}

% ============================================================================
\section{Prioriser les Vulnérabilités}
\label{sec:c21-prioriser}
% ============================================================================

Une analyse de vulnérabilités sur un grand parc produit des centaines de
résultats. On les trie en croisant la gravité technique (CVSS,
Chapitre~\ref{chap:menaces}) avec l'exposition, la valeur de l'actif et
l'exploitation connue.

\begin{tableaucours}{L{4.4cm}ccL{2.8cm}}{%
    \tableauHeader{Vulnérabilité} &
    \tableauHeader{CVSS} &
    \tableauHeader{Exposée} &
    \tableauHeader{Priorité}
}
Bibliothèque de documents du portail & 9,8 & Oui & 1 (72~h) \\
Version du serveur affichée & 5,3 & Oui & 3 \\
Faille d'un poste de test isolé & 8,8 & Non & 3 \\
Protocole SMBv1 actif sur un serveur & 8,1 & Interne & 2 (15~jours) \\
\end{tableaucours}
\vspace{-0.8em}
\begin{center}{\small\itshape Tableau~21.3. Prioriser selon la gravité et le contexte}\end{center}

\begin{boxalert}
\textbf{Violation typique.} Conclure « le système est sûr » parce que le
test n'a rien trouvé.

\smallskip
\textbf{Règle~:} un rapport dit « aucune vulnérabilité exploitable n'a
été trouvée dans le périmètre et la durée du test », jamais « le système
est sûr ». Ce qui n'a pas été testé reste inconnu.
\end{boxalert}

% ============================================================================
\section{Le Rapport}
\label{sec:c21-rapport}
% ============================================================================

\begin{tableaucours}{L{3.6cm}L{8.6cm}}{%
    \tableauHeader{Partie} &
    \tableauHeader{Contenu}
}
Synthèse & Pour la direction~: objectifs, résultat global, risques
principaux, en langage métier \\
Cadre & Périmètre, dates, méthodes, limites \\
Constatations & Pour chaque vulnérabilité~: description, preuve,
gravité, systèmes touchés \\
Recommandations & Correction, priorité, mesure provisoire \\
Annexes & Preuves, journaux du testeur, outils et versions \\
\end{tableaucours}
\vspace{-0.8em}
\begin{center}{\small\itshape Tableau~21.4. La structure d'un rapport de test}\end{center}

\begin{boxmetier}{Auditeur de sécurité}
\textbf{Faire le lien avec le risque.} L'auditeur traduit chaque
constatation en risque pour l'organisation et la rattache au registre
(Chapitre~\ref{chap:risques})~; il vérifie ensuite, par un test de
contrôle, que la correction est effective.
\end{boxmetier}

\begin{boxmetier}{RSSI}
\textbf{Commander et suivre.} Le RSSI choisit le type d'évaluation selon
la question posée, signe ou fait signer l'autorisation, et suit les
corrections jusqu'à leur vérification.
\end{boxmetier}

\separateur

\begin{boxresume}
\textbf{Ce qu'il faut retenir de ce chapitre~:}
\begin{itemize}
    \item \textbf{Analyse de vulnérabilités, test d'intrusion, audit et
          équipe rouge} répondent à des questions différentes.
    \item \textbf{Aucun test sans autorisation écrite}~: périmètre,
          méthodes, fenêtre, contacts.
    \item \textbf{Un test d'intrusion} planifie, découvre, attaque et rend
          compte~; en boîte noire, grise ou blanche.
    \item \textbf{On priorise} en croisant gravité, exposition, valeur et
          exploitation connue.
    \item \textbf{Un rapport} sépare synthèse, cadre, constatations et
          recommandations, et dit ses limites.
\end{itemize}
\end{boxresume}

\begin{boxinfo}
\textbf{Points de contrôle de l'élève.} Je sais~: choisir un type
d'évaluation~; rédiger une note d'autorisation~; nommer les quatre phases
d'un test~; prioriser des vulnérabilités~; structurer un rapport et
formuler ses limites.
\end{boxinfo}

\bigskip

\begin{boxexo}[title={\ding{45}\quad Exercice~21.1. Quel type~? \niveaubase}]
Quel type d'évaluation pour~: « nos serveurs ont-ils des correctifs
manquants~? »~; « un attaquant peut-il atteindre la paie depuis
Internet~? »~; « appliquons-nous notre politique d'accès~? »
\end{boxexo}

\begin{boxexo}[title={\ding{45}\quad Exercice~21.2. Les boîtes \niveaubase}]
Pour un test du portail de suivi de LiZbiZ, quelles informations le
testeur reçoit-il en boîte noire, grise et blanche~?
\end{boxexo}

\begin{boxexo}[title={\ding{45}\quad Exercice~21.3. La note \niveaumoyen}]
Rédigez la note d'autorisation d'un test d'intrusion du réseau de
l'agence de Garoua.
\end{boxexo}

\begin{boxexo}[title={\ding{45}\quad Exercice~21.4. Prioriser \niveaumoyen}]
Classez cinq vulnérabilités de votre choix, en justifiant chaque rang par
la gravité, l'exposition et la valeur de l'actif.
\end{boxexo}

\begin{boxexo}[title={\ding{45}\quad Exercice~21.5. La limite \niveaumoyen}]
Réécrivez la phrase « le portail de l'ANPP est sécurisé » d'un rapport,
de façon exacte.
\end{boxexo}

\begin{boxexo}[title={\ding{45}\quad Exercice~21.6. Le prestataire testeur \niveaumoyen}]
LiZbiZ veut confier le test d'intrusion à SAHEL INFOGÉRANCE, qui exploite
déjà ses centres. Qu'en pensez-vous~?
\end{boxexo}

\begin{boxexo}[title={\ding{45}\quad Exercice~21.7. Vérifier la correction \niveaumoyen}]
Comment vérifier que la bibliothèque vulnérable du portail a bien été
mise à jour, et que la vulnérabilité n'est plus exploitable~?
\end{boxexo}

\begin{boxexo}[title={\ding{45}\quad Exercice~21.8. Synthèse \niveaumoyen}]
Rédigez la synthèse d'une demi-page du test de l'ANPP pour son directeur
général.
\end{boxexo}

\begin{boxexo}[title={\ding{45}\quad Exercice~21.9. Un programme d'évaluation \niveauexpert}]
Proposez à la BANQUE SANAGA un programme annuel d'évaluation~: types,
périmètres, fréquences, prestataires, budget relatif.
\end{boxexo}

\begin{boxexo}[title={\ding{45}\quad Exercice~21.10. Rejouer GAROUA \niveauexpert}]
Concevez un exercice d'équipe rouge qui rejoue l'Opération GAROUA sur
l'architecture cible (Chapitre~\ref{chap:architecture})~: objectifs,
règles, mesures de la détection.
\end{boxexo}

\bigskip

\begin{boxinfo}
\textbf{Pour aller plus loin.}
\begin{itemize}
    \item NIST SP~800-115~\cite{nist2008sp800115}.
    \item ISO~19011 et ISO/IEC~27007~\cite{iso19011_2018,iso27007_2020}.
\end{itemize}
\end{boxinfo}

\begin{boxinfo}
\textbf{Transition.} L'évaluation est ponctuelle~; la surveillance est
permanente. Le Chapitre~\ref{chap:operations} traite des opérations de
sécurité et de la réponse aux incidents.
\end{boxinfo}
